\documentclass[11pt]{article}

\usepackage[a4paper,margin=27mm]{geometry}
\usepackage{amsmath,amssymb,amsthm,mathtools}
\usepackage{booktabs,longtable,array}
\usepackage[hidelinks]{hyperref}
\usepackage{microtype}
\usepackage{xcolor}
\usepackage{enumitem}

\definecolor{wrehblue}{RGB}{24,86,128}

\newtheorem{definition}{Definition}[section]
\newtheorem{proposition}[definition]{Proposition}
\newtheorem{example}[definition]{Example}
\newtheorem{remark}[definition]{Remark}

\newcommand{\W}{\mathcal W}
\newcommand{\Pcal}{\mathcal P}
\newcommand{\Acal}{\mathcal A}
\newcommand{\Bcal}{\mathcal B}
\newcommand{\Qcal}{\mathcal Q}
\newcommand{\QA}{\mathcal Q_{\mathcal A}}
\newcommand{\RA}{R_{\mathcal A}}

\title{\textbf{WR-I Preprint Candidate v0.3}\\[4pt]
\large Response-Defined World Spaces:\\
Finite-Resolution Equivalence and Identifiability of Global Realizations}
\author{A. A. Malachevsky\\
ORCID: 0009-0008-6009-3196\\
World Realizability \& Epistemic Horizons (WREH)\\
Research Programme \& Community}
\date{7 October 2026}

\begin{document}
\maketitle

\begin{abstract}
We formulate a response-based framework for studying admissible global
realizations that are indistinguishable under a declared family of observation
protocols. A world-response structure consists of a class of admissible global
realizations, a family of admissible protocols, and the corresponding response
maps. Exact response equality induces a quotient world space. Identifiable
quantities are precisely those that factor through this quotient, protocol
refinement induces canonical quotient maps, and a gauge-completeness criterion
separates representational redundancy from residual response
non-identifiability. For metric gauge quotients we prove a localized
finite-resolution separation result: on any compact admissible domain, joint
point separation by a continuous protocol family implies finite protocol
separation with a uniform positive response margin at every fixed nonzero
resolution. We also show why ordinary finite-tolerance indistinguishability is
not generally an equivalence relation. Examples illustrate finite refinement,
gauge reduction, infinite-dimensional response fibres, and the dependence of
identifiability on admissibility assumptions. The framework is ontologically
neutral. It does not identify response equivalence with physical identity, and
it makes no claim that observation creates, selects, or modifies a world.
\end{abstract}

\section{Introduction}

Inverse problems rarely begin with an unconstrained question of what reality
``is.'' They begin with a declared model class, a declared mode of access, and
a map from candidate structures to observable data. When that map is not
injective, the data determine an equivalence class rather than a unique hidden
object. This elementary fact appears in mature theories of statistical and
structural identifiability, in inverse boundary problems, and in results on
observationally indistinguishable spacetimes
\cite{Rothenberg1971,Saccomani2003,Bellu2007,Villaverde2016,MeshkatSullivant2014,Manchak2009}.

WR-I asks for a framework in which the compared objects are not merely
parameters or coefficients but \emph{admissible global realizations}. The word
``world'' is used only as shorthand for an element of such a declared class.
The basic question is:
\begin{quote}
Given a class of admissible global realizations and a family of physically
admissible protocols, which distinctions between realizations are actually
resolved by their responses?
\end{quote}
The answer is encoded by a response quotient. The quotient itself is standard
mathematics; the purpose of WR-I is to make it the controlled interface between
three layers that are often conflated: exact response identifiability, gauge
redundancy, and finite-resolution distinguishability.

\subsection*{Methodological precursor: from BC-XI to WREH}

The response-first viewpoint used here has a direct methodological precursor in
\emph{Boundary Compensation XI: The Inverse Isotypic Gap Problem and
Finite-Resolution Response Equivalence Classes} (BC-XI)
\cite{MalachevskyBCXI}. In BC-XI, a finite-resolution inverse response map could
be non-injective on hidden operator structures, so the inverse object was
naturally a response-equivalence class rather than a uniquely reconstructed
hidden representative.

WR-I retains that methodological lesson but changes the mathematical carrier.
The objects being compared are no longer hidden operators or finite-dimensional
realizations inside Boundary Compensation. They are admissible global
realizations $W\in\W$, compared through a declared family of admissible
protocols. Schematically,
\[
\begin{aligned}
\text{hidden structures}
&\xrightarrow{\text{BC-XI}}
\text{finite-resolution response classes},\\
\text{response classes}
&\xrightarrow{\text{WR-I}}
\text{classes of admissible global realizations}.
\end{aligned}
\]
This is a bridge, not an identification of programmes. No BC-XI theorem is
silently promoted into a statement about global worlds, spacetime, cosmology,
or ontology.

The distinction can be summarized in two questions:
\begin{quote}
BC-XI asks what a finite-resolution response determines about a hidden object.
WR-I asks what a family of admissible responses determines about a possible
global realization.
\end{quote}

\subsection*{Contributions and claim ceiling}

The paper makes five controlled framework contributions.
\begin{enumerate}[label=(\roman*)]
\item It defines a world-response structure $(\W,\Pcal,\{R_P\})$ and the exact
response quotient generated by a protocol bundle.
\item It characterizes identifiable quantities by factorization through that
quotient and organizes protocol refinement by canonical surjections between
quotients.
\item It separates gauge redundancy from residual response
non-identifiability and gives an exact gauge-completeness criterion.
\item It proves a localized finite-resolution separation proposition on
compact admissible domains, avoiding any global compactness assumption on the
entire physical state space.
\item It separates exact response equivalence from finite-tolerance
indistinguishability, which is generally non-transitive.
\end{enumerate}

The elementary factorization statements are not presented as new mathematical
theorems. The finite-resolution proposition is a compactness consequence
proved here for the WREH architecture; no historical priority claim is made for
that abstract topological statement. Recent domain-specific work also combines
exact observational quotients, monotone refinement, and explicitly
non-transitive finite-resolution ambiguity \cite{SimonsWashburn2026}. The
proposed novelty of WR-I is therefore architectural and programme-level, not a
claim to have invented observational equivalence, quotient spaces, structural
identifiability, gauge ambiguity, refinement of exact observational classes, or
the non-transitivity of finite-resolution ambiguity.

WR-I does \emph{not} claim that response-equivalent realizations are
ontologically identical, that underlying reality does not exist, that
measurement creates a world, that experiment rewrites the past, that the
Universe branches, or that a response quotient is a spacetime or cosmological
object.

\section{Relation to prior work}

\subsection{Structural identifiability}

Rothenberg's classical formulation of identification studies whether distinct
underlying structures induce the same observable probability law
\cite{Rothenberg1971}. In dynamical systems, structural identifiability asks
whether unknown parameters are uniquely recoverable from ideal input-output
data. Differential-algebraic approaches, global identifiability tests, and
identifiable reparametrizations form a mature literature
\cite{Saccomani2003,Bellu2007,Villaverde2016,MeshkatSullivant2014}.
WR-I therefore makes no novelty claim for input-output indistinguishability,
recoverable parameter combinations, or reparametrization as such.

A particularly close recent neighbour is Zhang's \emph{Mathematical
Identifiability Geometry}, which explicitly treats structures with identical
observation data as equivalence classes and studies geometric invariants of the
resulting space \cite{Zhang2026}. Simons and Washburn's 2026 spectroscopy
preprint gives an even closer domain-specific control: equality of observation
laws defines exact quotient classes, added modalities refine those classes,
and finite-resolution ambiguity is explicitly non-transitive
\cite{SimonsWashburn2026}. Consequently, WR-I does not claim novelty for any
of those ingredients individually. Its narrower target is the package
consisting of an admissible global-realization carrier, an explicit admissible
protocol family, separation of gauge orbits from residual response fibres, a
compact-domain finite-separation statement in a declared state-space metric,
and a handoff to a distinct global-realizability problem.

\subsection{Gauge ambiguity in inverse problems}

Inverse boundary problems provide concrete examples in which the observation
map is invariant under a known transformation group. In the anisotropic
Calder\'on problem, boundary measurements determine the interior geometry only
up to boundary-fixing changes of coordinates in the appropriate setting
\cite{LassasLiimatainenSalo2020}. This is a domain-specific benchmark for the
difference between representational redundancy and additional response
non-identifiability.

\subsection{Global spacetime underdetermination}

Manchak showed that global spacetime structure in classical general relativity
can remain observationally underdetermined in a robust sense
\cite{Manchak2009}. Cinti and Fano later argued that the physical
reasonableness of candidate indistinguishable spacetimes matters to the force
of such conclusions \cite{CintiFano2021}. For WR-I this is not a dispute to be
resolved but a structural lesson: the admissible realization class must be
part of the mathematical input. One cannot infer a physically meaningful
non-identifiability result from an unrestricted class of formally constructible
objects.

\subsection{Global-section obstructions}

Abramsky and Brandenburger formulate contextuality and non-locality through
obstructions to global sections of compatible empirical data
\cite{AbramskyBrandenburger2011}. Their work is a major predecessor for the
future WR-II question of whether locally or finitely compatible information
admits a global realization. WR-I does not identify a global section of a
measurement scenario with a global world realization; that bridge, if useful,
must be formulated explicitly.

\section{World-response structures}

\begin{definition}[World-response structure]
A \emph{world-response structure} is a triple
\[
\mathbf W=(\W,\Pcal,\{R_P\}_{P\in\Pcal}),
\]
where $\W$ is a declared class of \emph{admissible global realizations},
$\Pcal$ is a family of admissible observation protocols, and each protocol
$P$ has a response map
\[
R_P:\W\to Y_P
\]
into a protocol-dependent response space $Y_P$.
\end{definition}

\begin{remark}[Admissibility is part of the model]
The set $\W$ is not the class of all mathematically describable objects, and
$\Pcal$ is not the class of all mathematically definable readouts. Both are
inputs whose physical or modelling justification must be supplied by an
application.
\end{remark}

\begin{definition}[Protocol bundle and joint response]
For $\Acal\subseteq\Pcal$, define
\[
Y_{\Acal}:=\prod_{P\in\Acal}Y_P,
\qquad
R_{\Acal}(W):=(R_P(W))_{P\in\Acal}.
\]
For the empty bundle, $Y_{\varnothing}$ is the one-point set.
\end{definition}

\begin{definition}[Exact response equivalence]
For a protocol bundle $\Acal$, define
\[
W_1\sim_{\Acal}W_2
\quad\Longleftrightarrow\quad
R_{\Acal}(W_1)=R_{\Acal}(W_2).
\]
\end{definition}

\begin{definition}[Response-defined world space]
The response-defined world space for $\Acal$ is
\[
\QA:=\W/\!\sim_{\Acal},
\]
with quotient map
\[
q_{\Acal}:\W\to\QA,
\qquad
q_{\Acal}(W)=[W]_{\Acal}.
\]
\end{definition}

\section{Exact response quotient and identifiable quantities}

\begin{proposition}[Response quotient representation]\label{prop:quotient}
The relation $\sim_{\Acal}$ is an equivalence relation, and there is a unique
bijection
\[
\overline R_{\Acal}:\QA\longrightarrow\operatorname{Im}R_{\Acal}
\]
such that
\[
R_{\Acal}=\overline R_{\Acal}\circ q_{\Acal}.
\]
\end{proposition}

\begin{proof}
The relation is the kernel equivalence relation of $R_{\Acal}$. Define
$\overline R_{\Acal}([W]_{\Acal})=R_{\Acal}(W)$. Equality of classes is exactly
what makes this definition well posed and injective; surjectivity onto the
image and uniqueness follow from the definitions.
\end{proof}

\begin{remark}
As a bare set, $\QA$ contains exactly the distinctions present in the joint
response image. A topology, metric, smooth structure, measure, stratification,
or physical interpretation on $\QA$ requires additional assumptions.
\end{remark}

\begin{definition}[Identifiable quantity]
A map $F:\W\to Z$ is \emph{$\Acal$-identifiable} if
\[
W_1\sim_{\Acal}W_2
\quad\Longrightarrow\quad
F(W_1)=F(W_2).
\]
\end{definition}

\begin{proposition}[Identifiability factorization]\label{prop:identifiable}
A map $F:\W\to Z$ is $\Acal$-identifiable if and only if there exists a unique
map
\[
\overline F:\QA\to Z
\]
such that
\[
F=\overline F\circ q_{\Acal}.
\]
\end{proposition}

\begin{proof}
An $\Acal$-identifiable map is constant on the fibres of $q_{\Acal}$, so
$\overline F([W]_{\Acal})=F(W)$ is well defined. The converse is immediate from
the factorization, and uniqueness follows from surjectivity of $q_{\Acal}$.
\end{proof}

\section{Protocol refinement}

\begin{definition}[Refinement order]
For protocol bundles $\Acal,\Bcal\subseteq\Pcal$, write
\[
\Acal\preceq\Bcal
\quad\Longleftrightarrow\quad
\Acal\subseteq\Bcal.
\]
\end{definition}

\begin{proposition}[Refinement monotonicity]\label{prop:refinement}
If $\Acal\subseteq\Bcal$, then
\[
[W]_{\Bcal}\subseteq[W]_{\Acal}
\qquad\text{for every }W\in\W.
\]
Moreover, there is a unique surjection
\[
\pi_{\Bcal\Acal}:\mathcal Q_{\Bcal}\twoheadrightarrow\mathcal Q_{\Acal},
\qquad
[W]_{\Bcal}\mapsto[W]_{\Acal},
\]
satisfying
\[
q_{\Acal}=\pi_{\Bcal\Acal}\circ q_{\Bcal}.
\]
For $\Acal\subseteq\Bcal\subseteq\mathcal C$,
\[
\pi_{\mathcal C\Acal}
=
\pi_{\Bcal\Acal}\circ\pi_{\mathcal C\Bcal}.
\]
\end{proposition}

\begin{proof}
Equality of all responses in $\Bcal$ implies equality of the subfamily indexed
by $\Acal$. The stated quotient map is therefore well defined and surjective;
the commuting and composition identities follow directly.
\end{proof}

\begin{remark}
The directed hierarchy of response quotients is the interface to WR-II. WR-I
does not yet infer global existence from a refinement chain and does not assign
physical time or dynamics to the refinement parameter.
\end{remark}

\section{Gauge redundancy and residual non-identifiability}

\begin{definition}[Response-invariant gauge action]
Let a group $G$ act on $\W$. The action is response-invariant for $\Acal$ if
\[
R_P(g\cdot W)=R_P(W)
\]
for every $g\in G$, $W\in\W$, and $P\in\Acal$.
\end{definition}

\begin{proposition}[Gauge factorization]\label{prop:gauge-factor}
If the $G$-action is response-invariant for $\Acal$, then
\[
G\cdot W\subseteq[W]_{\Acal}
\]
for every $W\in\W$, and $R_{\Acal}$ factors uniquely through the orbit quotient:
\[
\W\xrightarrow{\pi_G}\W/G
\xrightarrow{\widetilde R_{\Acal}}Y_{\Acal}.
\]
\end{proposition}

\begin{proof}
Response invariance gives $R_{\Acal}(g\cdot W)=R_{\Acal}(W)$, so each orbit lies
inside one response class. Thus
$\widetilde R_{\Acal}(G\cdot W):=R_{\Acal}(W)$ is well defined and gives the
required factorization.
\end{proof}

\begin{proposition}[Gauge-completeness criterion]\label{prop:gauge-complete}
Assume the $G$-action is response-invariant for $\Acal$. The following are
equivalent:
\begin{enumerate}[label=(\alph*)]
\item $[W]_{\Acal}=G\cdot W$ for every $W\in\W$;
\item the induced map
\[
\widetilde R_{\Acal}:\W/G\to Y_{\Acal}
\]
is injective.
\end{enumerate}
\end{proposition}

\begin{proof}
If every response class is exactly one orbit, equal induced responses imply
equal orbits. Conversely, Proposition~\ref{prop:gauge-factor} gives
$G\cdot W\subseteq[W]_{\Acal}$; injectivity on $\W/G$ forces any
$W'\in[W]_{\Acal}$ to belong to the same orbit as $W$.
\end{proof}

\begin{remark}
When Proposition~\ref{prop:gauge-complete}(b) fails, the observation system
leaves distinct gauge orbits unresolved. WR-I calls this \emph{residual response
non-identifiability}; no ontological conclusion follows from that label alone.
\end{remark}

\section{Localized finite-resolution separation}

Global compactness of $X=\W/G$ is unnecessarily strong for the finite-resolution
question. Experimental design normally concerns a bounded or otherwise
controlled admissible domain. We therefore localize the argument.

Assume that the full protocol family $\Pcal$ is $G$-invariant, so every
response descends to the orbit quotient. Assume further that $X=\W/G$ is a
metric space with metric $d_X$. Let $K\subseteq X$ be a compact admissible
domain. Suppose each descended protocol response is continuous,
\[
\widetilde R_P:X\to(Y_P,d_P),
\]
and that the full admissible family separates points of $K$:
\[
x\neq y,
\quad x,y\in K
\quad\Longrightarrow\quad
\exists P\in\Pcal:
\widetilde R_P(x)\neq\widetilde R_P(y).
\]

\begin{proposition}[Localized finite-resolution separation]\label{prop:finite-resolution}
For every $\delta>0$ there exists a finite protocol bundle
\[
\Acal_{K,\delta}=\{P_1,\dots,P_m\}\subseteq\Pcal
\]
and, whenever
\[
K^{(2)}_{\delta}
:=
\{(x,y)\in K\times K:d_X(x,y)\ge\delta\}
\]
is nonempty, a constant $\eta_{K,\delta}>0$ such that
\[
x,y\in K,
\qquad
d_X(x,y)\ge\delta
\]
implies
\[
\max_{1\le j\le m}
 d_{P_j}\!\left(
 \widetilde R_{P_j}(x),
 \widetilde R_{P_j}(y)
 \right)
\ge\eta_{K,\delta}.
\]
If $K^{(2)}_{\delta}$ is empty, the separation statement is vacuous.
\end{proposition}

\begin{proof}
The set $K^{(2)}_{\delta}$ is closed in the compact space $K\times K$, hence
compact. For every $P\in\Pcal$ define
\[
U_P
=
\{(x,y)\in K^{(2)}_{\delta}:
 d_P(\widetilde R_P(x),\widetilde R_P(y))>0\}.
\]
Continuity makes $U_P$ open in $K^{(2)}_{\delta}$. Point separation implies
that the family $\{U_P\}_{P\in\Pcal}$ covers $K^{(2)}_{\delta}$. Compactness
therefore yields a finite subcover $U_{P_1},\dots,U_{P_m}$.

Define
\[
h(x,y)
=
\max_{1\le j\le m}
 d_{P_j}(\widetilde R_{P_j}(x),\widetilde R_{P_j}(y)).
\]
The function $h$ is continuous and strictly positive on
$K^{(2)}_{\delta}$. Hence it attains a positive minimum,
\[
\eta_{K,\delta}
:=
\min_{(x,y)\in K^{(2)}_{\delta}}h(x,y)>0.
\]
\end{proof}

\begin{remark}[What compactness is actually needed]
The proof only requires compactness of the separated-pair set
$K^{(2)}_{\delta}$, not compactness of the whole physical state space $X$.
Local compactness of $X$ alone does not produce the finite subcover needed for
a global statement; one must still select a compact domain or otherwise prove
compactness of the relevant pair set.
\end{remark}

\begin{remark}[Resolution dependence]
The proposition is not a finite exact-reconstruction theorem. As
$\delta\downarrow0$, the required bundle size may grow and
$\eta_{K,\delta}$ may tend to zero. The chosen metric, compact domain,
continuity assumptions, and protocol admissibility are all application-level
inputs.
\end{remark}

\section{Finite tolerance is not an equivalence relation in general}

Suppose each response space $Y_P$ carries a metric $d_P$, and choose positive
weights $w_P$ for which
\[
d_{\Acal}(W_1,W_2)
=
\sup_{P\in\Acal}w_P d_P(R_P(W_1),R_P(W_2))
\]
is finite on the pairs under consideration. For $\varepsilon>0$, define
\[
W_1\approx_{\Acal,\varepsilon}W_2
\quad\Longleftrightarrow\quad
 d_{\Acal}(W_1,W_2)\le\varepsilon.
\]

\begin{proposition}[Failure of transitivity]\label{prop:tolerance}
The relation $\approx_{\Acal,\varepsilon}$ is generally not transitive and
therefore does not generally define a quotient.
\end{proposition}

\begin{proof}
For one scalar response choose
\[
R(W_0)=0,
\qquad
R(W_1)=\frac34\varepsilon,
\qquad
R(W_2)=\frac32\varepsilon.
\]
Then $W_0\approx_\varepsilon W_1$ and
$W_1\approx_\varepsilon W_2$, but
$W_0\not\approx_\varepsilon W_2$.
\end{proof}

\begin{remark}
WR-I therefore reserves ``equivalence class'' for exact response equality
unless an explicit equivalence closure or another justified coarse relation is
introduced. Finite-precision objects are instead called tolerance fibres,
response neighbourhoods, or uncertainty regions.
\end{remark}

\section{Examples}

\begin{example}[Finite binary realization with gauge and hidden physical content]
Let
\[
\W=\{0,1\}^4,
\qquad
W=(a,b,g,h).
\]
Let $G=\mathbb Z_2$ flip the label $g$ and leave $(a,b,h)$ fixed. Define
\[
R_a(W)=a,
\qquad
R_b(W)=b,
\qquad
R_h(W)=a\oplus h.
\]
With only $R_a$, there are two response classes of size eight. With
$(R_a,R_b)$, there are four classes of size four and the physical variable $h$
remains unresolved. With $(R_a,R_b,R_h)$, there are eight classes of size two,
and each class is exactly one gauge orbit. Thus the induced response map is
injective on $\W/G$.

An artificial readout $R_g(W)=g$ breaks gauge invariance and would make
representatives distinguishable. The example illustrates why admissible
protocols must be part of the model definition rather than an arbitrary set of
mathematical readouts.
\end{example}

\begin{example}[Smooth refinement from genuine ambiguity to pure gauge]
Let
\[
\W=S^1_x\times S^1_g,
\]
with gauge group $G=S^1$ acting by
\[
\alpha\cdot(x,g)=(x,g+\alpha).
\]
Then $\W/G\cong S^1_x$. Consider
\[
R_c(x,g)=\cos x,
\qquad
R_s(x,g)=\sin x.
\]
For $\Acal=\{R_c\}$, the induced map
$x\mapsto\cos x$ is non-injective, so generic response fibres contain two
distinct gauge orbits. For
$\Bcal=\{R_c,R_s\}$, the induced map
\[
x\mapsto(\cos x,\sin x)
\]
is injective, and the remaining response classes are exactly the gauge orbits.
Thus protocol refinement removes residual non-identifiability without removing
gauge redundancy.
\end{example}

\begin{example}[Continuous histories, dense observations, and admissibility]
Let
\[
\W=C([0,1],\mathbb R),
\]
and for each $t\in[0,1]$ let $P_t$ denote the evaluation protocol with
\[
R_{P_t}(f)=f(t).
\]
For a finite sampling bundle $\Acal=\{P_{t_1},\ldots,P_{t_n}\}$,
\[
[f]_{\Acal}
=
\{g\in C([0,1]):g(t_i)=f(t_i),\ i=1,\ldots,n\}
=
f+V_{\Acal},
\]
where
\[
V_{\Acal}
=
\{u\in C([0,1]):u(t_i)=0\text{ for all }i\}.
\]
The space $V_{\Acal}$ is infinite-dimensional, so every finite sampling bundle
leaves an infinite-dimensional response fibre.

Let $D\subset[0,1]$ be countable and dense. If two continuous histories agree
on every point of $D$, continuity forces them to agree everywhere. Thus the countable protocol family $\{P_t:t\in D\}$ is jointly injective on
this admissible realization class.

The conclusion fails after enlarging the admissible class. Choose
$t_0\in[0,1]\setminus D$, let
\[
f(t)=0,
\]
and define the bounded discontinuous history
\[
g(t)=
\begin{cases}
1,&t=t_0,\\
0,&t\neq t_0.
\end{cases}
\]
Then $f$ and $g$ agree on every protocol in the dense family $D$ but are not the
same global realization. Thus the same protocol family can be identifying on
one admissible class and non-identifying on a larger one.
\end{example}

\section{Calder\'on problem as a gauge benchmark}

The anisotropic Calder\'on problem supplies a concrete benchmark for
Propositions~\ref{prop:gauge-factor}--\ref{prop:gauge-complete}. Boundary
voltage-current data are invariant under the appropriate boundary-fixing
coordinate transformations, and uniqueness results are therefore formulated
up to that invariance \cite{LassasLiimatainenSalo2020}. In the language of
WR-I, the known transformation orbit lies inside the response fibre; the
substantive inverse problem is whether any additional ambiguity remains after
passing to the quotient by that orbit.

This analogy is deliberately limited. WR-I does not claim that every gauge
problem is a Calder\'on problem or that the abstract group $G$ has a universal
physical meaning. The example serves only to calibrate the distinction between
representational redundancy and residual non-identifiability.

\section{Prior-art and novelty ledger}

\begin{longtable}{>{\raggedright\arraybackslash}p{0.20\textwidth}
                  >{\raggedright\arraybackslash}p{0.36\textwidth}
                  >{\raggedright\arraybackslash}p{0.34\textwidth}}
\toprule
\textbf{Prior line} & \textbf{Established idea} & \textbf{WR-I boundary} \\
\midrule
\endfirsthead
\toprule
\textbf{Prior line} & \textbf{Established idea} & \textbf{WR-I boundary} \\
\midrule
\endhead

Rothenberg (1971) &
Identification and observational equivalence in parametric stochastic models. &
No novelty claim for observational equivalence or identification as such. \\[5pt]

Saccomani et al. (2003); Bellu et al. (2007); Villaverde et al. (2016) &
Structural identifiability of dynamical models from ideal input-output data,
including differential-algebraic and computational tests. &
WR-I does not relabel standard parameter identifiability as world
identifiability; it changes the carrier to an explicitly declared class of
global realizations and an explicit protocol family. \\[5pt]

Meshkat--Sullivant (2014) &
Identifiable reparametrizations and parameter combinations using algebraic and
graph-theoretic structure. &
Factorization through identifiable combinations is prior art; WR-I uses a
quotient language at the global-realization level. \\[5pt]

Zhang (2026) &
Geometric treatment of observation-equivalence classes, including
indistinguishability strata and unique-identification points. &
WR-I does not claim that geometry of observation classes is new. \\[5pt]

Simons--Washburn (2026) &
A physical measurement-design framework with exact observation-law quotients,
monotone refinement under added modalities, and explicitly non-transitive
finite-resolution ambiguity. &
These ingredients are prior art in a domain-specific finite-resolution setting.
WR-I's remaining distinction is the global-realization carrier, explicit gauge
separation, compact-domain finite separation in a state-space metric, and the
WR-II global-realizability handoff. \\[5pt]

BC-XI (2026) &
Finite-resolution response equivalence as the inverse object when hidden
operator structures are not uniquely reconstructed. &
Methodological and historical precursor. WR-I changes the carrier from hidden
operator structures to admissible global realizations; no physical or
ontological BC claim is imported. \\[5pt]

Manchak (2009) &
Observational underdetermination of global spacetime structure in classical
GR. &
Domain-specific predecessor for global-world underdetermination; WR-I abstracts
the observation mechanism into a declared protocol family. \\[5pt]

Cinti--Fano (2021) &
Critique stressing physical reasonableness of observationally indistinguishable
spacetimes. &
Supports treating $\W_{\mathrm{admissible}}$ as part of the formal input. \\[5pt]

Abramsky--Brandenburger (2011) &
Contextuality/non-locality as obstruction to global sections of local empirical
data. &
Major predecessor for WR-II. WR-I does not equate global sections with global
world realizations. \\[5pt]

Anisotropic Calder\'on inverse problem &
Boundary response invariance under boundary-fixing coordinate changes and
uniqueness up to the corresponding gauge in suitable settings. &
Benchmark for gauge factorization and the question whether response fibres are
strictly larger than gauge orbits. \\
\bottomrule
\end{longtable}

The defensible novelty claim at v0.3 is therefore architectural:
\begin{quote}
WR-I formulates a protocol-indexed framework whose primary carrier is a
declared class of admissible global realizations, explicitly distinguishes
known gauge orbits from residual response fibres, and proves that joint
point-separation on a compact admissible domain yields a finite protocol family
with a uniform positive response margin for every fixed nonzero separation in
the declared state-space metric. These objects are organized as the input
language for a distinct global-realizability problem rather than as a claim to
a preferred ontology.
\end{quote}

The claim remains subject to continuing literature audit. In particular, no
priority claim is attached to the elementary quotient factorization statements
or to the general use of compactness in finite subcover arguments.

\section{Discussion: three levels that must remain distinct}

The framework separates three questions.

\paragraph{Exact response identifiability.}
For a fixed protocol bundle, exact equality of responses produces a genuine
equivalence relation and quotient. This is an idealized structural question.

\paragraph{Gauge-reduced identifiability.}
If a declared gauge group acts response-invariantly, one must first determine
whether the observation system distinguishes gauge orbits. Failure of
injectivity on $\W/G$ is stronger than mere representational redundancy.

\paragraph{Finite-resolution distinguishability.}
Real experiments have finite accuracy and operate on controlled domains. The
localized finite-resolution result therefore concerns a compact admissible
region $K$ and a nonzero resolution scale $\delta$, not a global exact
reconstruction claim.

These distinctions are also a claim firewall. Neither persistent response
multiplicity nor finite-resolution ambiguity implies that multiple worlds are
physically instantiated. Conversely, a unique point of a declared quotient is
uniqueness relative to the chosen realization class and protocol family; it is
not automatically uniqueness of ontology.

\section{Handoff to WR-II}

WR-I ends where a new mathematical question begins. Suppose progressively richer
finite protocol bundles produce compatible response constraints. Even if every
finite stage is nonempty, one may ask whether there exists any admissible global
realization satisfying all stages simultaneously.

Schematically, refinement gives
\[
\mathcal Q_1
\longleftarrow
\mathcal Q_2
\longleftarrow
\mathcal Q_3
\longleftarrow\cdots.
\]
WR-II will study the distinction between
\[
\text{finite compatibility}
\qquad\text{and}\qquad
\text{global realizability}.
\]
Potential tools include compactness principles, inverse/projective limits,
extension theorems, explicit counterexamples, and comparison with
sheaf-theoretic global-section obstructions. None of these is assumed by WR-I.
In particular, protocol refinement is not physical time, and failure of global
realizability is not interpreted as destruction or nonexistence of the physical
Universe without a separate bridge.

\section*{Conclusion}

WR-I provides a deliberately modest mathematical interface. Responses define
which distinctions survive a declared observation system; gauge structure
identifies a known class of redundant distinctions; compact-domain separation
relates ideal point separation to finite experimental design; and admissibility
controls what counts as a candidate global realization in the first place.

The intellectual route from BC-XI to WR-I is therefore continuous but not
identical: BC-XI established a response-first inverse viewpoint for hidden
finite-resolution structures, whereas WR-I lifts the question to classes of
admissible global realizations. The next step is not a stronger ontological
claim but the mathematical question of global completion.

\begin{thebibliography}{99}

\bibitem{Rothenberg1971}
T. J. Rothenberg,
\emph{Identification in Parametric Models},
Econometrica 39(3) (1971), 577--591.
DOI: 10.2307/1913267.

\bibitem{Saccomani2003}
M. P. Saccomani, S. Audoly, and L. D'Angi\`o,
\emph{Parameter identifiability of nonlinear systems: the role of initial conditions},
Automatica 39(4) (2003), 619--632.
DOI: 10.1016/S0005-1098(02)00302-3.

\bibitem{Bellu2007}
G. Bellu, M. P. Saccomani, S. Audoly, and L. D'Angi\`o,
\emph{DAISY: a new software tool to test global identifiability of biological
and physiological systems},
Computer Methods and Programs in Biomedicine 88(1) (2007), 52--61.
DOI: 10.1016/j.cmpb.2007.07.002.

\bibitem{Villaverde2016}
A. F. Villaverde, A. Barreiro, and A. Papachristodoulou,
\emph{Structural Identifiability of Dynamic Systems Biology Models},
PLOS Computational Biology 12(10) (2016), e1005153.
DOI: 10.1371/journal.pcbi.1005153.

\bibitem{MeshkatSullivant2014}
N. Meshkat and S. Sullivant,
\emph{Identifiable reparametrizations of linear compartment models},
Journal of Symbolic Computation 63 (2014), 46--67.
DOI: 10.1016/j.jsc.2013.11.002.

\bibitem{Zhang2026}
J. Zhang,
\emph{Mathematical Identifiability Geometry},
Zenodo preprint (2026).
DOI: 10.5281/zenodo.22023251.


\bibitem{SimonsWashburn2026}
M. Simons and J. Washburn,
\emph{Finite-Resolution Identifiability and Measurement Design for Molecular
Conformer Spectroscopy},
arXiv:2608.02637 (2026).

\bibitem{MalachevskyBCXI}
A. A. Malachevsky,
\emph{Boundary Compensation XI: The Inverse Isotypic Gap Problem and
Finite-Resolution Response Equivalence Classes},
Zenodo (2026).
DOI: 10.5281/zenodo.20748061.

\bibitem{Manchak2009}
J. B. Manchak,
\emph{Can we know the global structure of spacetime?},
Studies in History and Philosophy of Modern Physics 40(1) (2009), 53--56.
DOI: 10.1016/j.shpsb.2008.07.004.

\bibitem{CintiFano2021}
E. Cinti and V. Fano,
\emph{Careful with those scissors, Eugene! Against the observational
indistinguishability of spacetimes},
Studies in History and Philosophy of Science 89 (2021), 103--113.
DOI: 10.1016/j.shpsa.2021.07.007.

\bibitem{AbramskyBrandenburger2011}
S. Abramsky and A. Brandenburger,
\emph{The Sheaf-Theoretic Structure of Non-Locality and Contextuality},
New Journal of Physics 13 (2011), 113036.
DOI: 10.1088/1367-2630/13/11/113036.

\bibitem{LassasLiimatainenSalo2020}
M. Lassas, T. Liimatainen, and M. Salo,
\emph{The Poisson embedding approach to the Calder\'on problem},
Mathematische Annalen 377 (2020), 19--67.
DOI: 10.1007/s00208-019-01818-3.

\end{thebibliography}

\end{document}
